\begin{abstract}
\noindent
La segmentación automática de tumores cerebrales en resonancia magnética es un problema de visión por computador con un fuerte desequilibrio de clases: en el conjunto de datos utilizado, formado por 3064 cortes ponderados en T1 con contraste de 233 pacientes, el tumor ocupa de media el 1.7\,\% de los píxeles de cada imagen. Este trabajo aborda la segmentación multiclase de meningiomas, gliomas y tumores hipofisarios como un \emph{problema de optimización en dos niveles}. En el nivel interno, el entrenamiento de una red U-Net se formula como la minimización de una función de pérdida diferenciable y no convexa sobre entre $1.8$ y $7.2$ millones de parámetros, resuelta mediante descenso de gradiente estocástico adaptativo (Adam) con gradientes exactos obtenidos por retropropagación. Se analizan desde el punto de vista de la optimización la arquitectura, las funciones de pérdida sustitutas (Dice, Focal y Focal Tversky) y la programación de la tasa de aprendizaje. En el nivel externo, los hiperparámetros que condicionan ese proceso (tasa de aprendizaje, decaimiento de pesos, número de filtros y profundidad) se ajustan con un algoritmo genético de codificación real, con selección por torneo, cruce BLX-$\alpha$, mutación gaussiana con reflexión y elitismo, y una función de aptitud de baja fidelidad basada en el coeficiente Dice de validación tras tres épocas de entrenamiento. Con 46 evaluaciones, el algoritmo genético mejora la aptitud de la configuración manual de 0.441 a 0.509. Reentrenada a longitud completa, la configuración evolucionada eleva el Dice medio en el conjunto de prueba de 0.691 a 0.776 y alcanza el mejor valor de validación de la referencia en 10 épocas en lugar de 35. El análisis posterior de la solución heurística revela un sesgo de horizonte corto en la función de aptitud, que favorece redes pequeñas de convergencia rápida. Corregido mediante recocido coseno de la tasa de aprendizaje, mayor capacidad (32 filtros) y una calibración de los pesos de la pérdida Focal Tversky guiada por diagnóstico, el modelo entrenado desde cero alcanza un Dice medio de 0.8175 en prueba. Como complemento, un codificador ResNet34 preentrenado en ImageNet obtiene 0.8250, un resultado estadísticamente equivalente que resuelve la confusión entre tipos tumorales pero no mejora la delimitación del glioma. El algoritmo genético explica aproximadamente dos tercios de la mejora total respecto a la configuración manual.

\medskip
\noindent\textbf{Palabras clave:} optimización en dos niveles, descenso de gradiente, retropropagación, U-Net, algoritmo genético, optimización de hiperparámetros, pérdida Focal Tversky, segmentación de tumores cerebrales.
\end{abstract}
